\subsection{Adjoint systems and the strict genus inequality}
\label{subsec:canonical-subsystem-strictness}

A logarithmic form vanishing on a very ample curve controls much more
than its first nonzero coefficient. It gives a system of canonical
forms whose rational map has surface image. We first prove this
without a conductor equality. We then use the positive square to split an Euler extension at
conductor equality. The resulting structure theorem excludes the
first four conductors and gives the strict genus inequality.

\begin{proposition}\label{prop:residual-adjoint}
Let $X$ be a smooth projective integral surface over an algebraically
closed field $k$ of characteristic $p>0$, let $Y\subset X$ be a smooth
integral very ample divisor, and let $s$ be its canonical section.
Suppose that a nonzero closed global one-form $\omega$ vanishes along
$Y$. Let $m$ be its exact order along $Y$, and let $Z$ be the
divisorial zero part, with multiplicities, of
$\theta=\omega/s^m\in H^0(X,\Omega^1_{X/k}(-mY))$.
Then $m=pa$ for an integer $a\geq1$, $Y$ is not a component of $Z$, and
\[
 0\leq Z\cdot Y\leq 2g(Y)-2-mY^2.
\]
The divisor
\[
 A=K_X-(m-2)Y-Z
\]
is nef and big, with $A^2>0$.
More precisely, the linear system
\[
 W_K=\{d((u/s)\omega):u\in H^0(X,\OO_X(Y))\}
 \subset H^0(X,\omega_X)
\]
has exact fixed divisor $(m-2)Y+Z$ and two-dimensional image.
Its residual system has no fixed curve, and on a dense open of $Y$
its rational map is the given hyperplane embedding. The map
$u\mapsto d((u/s)\omega)$ has kernel $ks$. In particular,
\begin{align*}
 h^0(X,\Omega^1_{X/k})&\geq h^0(X,\OO_X(Y)),\\
 \operatorname{rank}\bigl[d:H^0(X,\Omega^1_{X/k})\to
 H^0(X,\omega_X)\bigr]&\geq h^0(X,\OO_X(Y))-1,\\
 p_g(X)&\geq h^0(X,\OO_X(Y))-1.
\end{align*}
Both $K_X-(m-2)Y$ and $K_X-(p-2)Y$ are big, and the canonical
map of $X$ has two-dimensional image.
\end{proposition}

\begin{proof}
The leading restriction of $\theta$ to $Y$ is nonzero. Its tangential
image is nonzero as well, because the kernel on global sections in
the twisted conormal sequence is
$H^0(Y,\OO_Y(-(m+1)Y))=0$. Locally write
$\omega=t^m(b\,dt+\eta)$, with $\eta$ nonzero tangentially modulo $t$.
The term of order $m-1$ in $d\omega$ is $m\,dt\wedge\eta$.
Thus $p\mid m$, and $m\geq p\geq2$.
The nonzero tangential section of $\omega_Y(-mY)$ gives
$2g(Y)-2\geq mY^2>0$, so $g(Y)\geq2$.

Regular local rings are factorial
\cite[Lemma~15.123.2, Tag~0AG0]{Stacks}.
The minimum valuations of the two coefficients of $\theta$ therefore
define an effective Cartier divisor $Z$. Dividing those coefficients
by their common prime factors gives a regular section
\[
 \theta'\in H^0(X,\Omega^1_{X/k}(-mY-Z)),\qquad
 \theta=s_Z\theta',
\]
with no divisorial zero. Exactness of $m$ excludes $Y$ from $Z$.
The conormal term after restriction to $Y$ has degree
$-(m+1)Y^2-Z\cdot Y<0$. Hence the tangential image
\[
 \sigma'\in H^0(Y,\omega_Y(-mY-Z))
\]
is nonzero. Its finite zero divisor has degree
$2g(Y)-2-mY^2-Z\cdot Y$, proving the first bound. Notice also that
\[
 A\cdot Y=Y^2+2g(Y)-2-mY^2-Z\cdot Y>0.
\]
Division by $s_Z$ need not preserve closedness; the argument below
uses closedness before that division.

\paragraph{The two-form system.}
Choose local frames $e_i$ of $\OO_X(Y)$ and write
$s=t_i e_i$, $u=u_i e_i$, and $\eta_i=t_i^{-m}\omega$.
The forms $\eta_i$ are regular and closed because $p\mid m$.
If $t_i=v_{ij}t_j$, then $u_i=v_{ij}u_j$ and
$\eta_i=v_{ij}^{-m}\eta_j$. The identity
\[
 t_i\,du_i-u_i\,dt_i
 =v_{ij}^2(t_j\,du_j-u_j\,dt_j)
\]
shows that
\begin{equation}\label{eq:residual-adjoint-local}
 \gamma_i(u)=(t_i\,du_i-u_i\,dt_i)\wedge\eta_i
\end{equation}
glues to a section of $K_X-(m-2)Y$.
Every such section is divisible by $Z$; denote the quotient by
$\gamma'(u)\in H^0(X,\OO_X(A))$ and its image space by $W$.

Along $Y$, equation~\eqref{eq:residual-adjoint-local} becomes
$-u_i\,dt_i\wedge\eta_i$. After dividing by $Z$, adjunction identifies
this restriction with $-u|_Y\,\sigma'$ in
$A|_Y=\omega_Y(-(m-1)Y-Z)$.
Multiplication by the nonzero section $\sigma'$ is injective on
sections of the integral curve. The exact sequence
\[
 0\longrightarrow\OO_X\xrightarrow{s}\OO_X(Y)
 \longrightarrow\OO_Y(Y)\longrightarrow0
\]
therefore gives $\ker\gamma'=ks$. Away from the zeros of $\sigma'$,
the rational map of $W$ is precisely the original embedding of $Y$
in the hyperplane defined by $s$. The restriction need not have its
complete linear series, and it may have finite base points.

On $X\setminus Y$, the functions $u/s$ are affine coordinates of the
very ample embedding and their differentials generate $\Omega^1_{X/k}$.
Up to invertible factors, the sections of $W$ are
$d(u/s)\wedge\theta'$. For a free rank-two module, the ideal of these
wedges is the ideal of the two coefficients of $\theta'$.
Thus no curve outside $Y$ is fixed. At the generic point of $Y$,
$\sigma'\ne0$, so $Y$ is not fixed either. This proves that $W$ has
no fixed curve and that the exact fixed divisor before division is $Z$.
It also proves nefness of $A$: for each integral curve $E$, choose a
member of $W$ not vanishing identically on $E$. Its effective zero
divisor has nonnegative intersection with $E$.

\paragraph{Why the system has surface image.}
Resolve its finite base locus by point blowups
\cite[Lemma~54.4.3, Tag~0C5H]{Stacks}, obtaining
\[
 \rho:R\longrightarrow X,\qquad
 \psi:R\longrightarrow\PP(W^\vee).
\]
The centres may lie on $Y$ or be infinitely near it.
The strict transform $Y'$ is isomorphic to $Y$: blowing up a smooth
curve at a point does not change that curve. For example, for local
parameters $t,x$ with $Y=(t=0)$, the chart $t=xv$ has strict transform
$v=0$, mapping isomorphically to $Y$.
The morphism $\psi|_{Y'}$ agrees with the original embedding on a dense
open. Since the target is separated, it agrees everywhere.

Suppose the reduced image $B=\psi(R)$ is a curve. It equals the
embedded image of $Y'$, so $B\simeq Y$ and $Y'$ is a section of
$\psi:R\to B$. Every exceptional curve of $\rho$ is a projective line
\cite[Section~54.3, Tag~0AGP]{Stacks}, and maps constantly to $B$.
To include inseparable maps in this assertion, a nonconstant map
$\PP^1\to B$ with nonzero differential would give a nonzero map from
a line bundle of positive degree to $\omega_{\PP^1}$, of degree $-2$.
If its differential is zero, its function-field image lies in
$k(t^p)$. Factoring through this rational function field and extending
the rational map by properness divides the degree by $p$.
Iteration eventually gives a nonzero differential and the same
contradiction. Here $g(B)=g(Y)\geq2$ is essential.

Keep the full Cartier pullback
\[
 \rho^*Y=Y'+\sum_i r_iE_i,\qquad r_i\geq0.
\]
All the exceptional components $E_i$ are vertical over $B$.
Set $F=k(B)$ and $\Gamma=R\times_B\Spec F$.
This is a proper integral regular curve: its local rings are
localizations of regular local rings of $R$.
Smoothness over the possibly imperfect field $F$ has not been assumed.
The section $Y'$ gives a single reduced $F$-rational point on $\Gamma$;
its equation is a uniformizer in the corresponding discrete valuation
ring. Therefore the globally generated line bundle
$\OO_R(\rho^*Y)|_\Gamma$ has degree one.

A nonzero section of this line bundle has a zero divisor consisting
of one reduced $F$-rational point. Global generation gives a second
section avoiding that point. The two sections define a finite morphism
$\Gamma\to\PP^1_F$ of degree one, by the degree formula on proper curves
\cite[Lemma~37.44.1, Tag~02LS; Lemma~33.44.11, Tag~0AYQ]{Stacks}.
Finite birationality onto a normal target makes this an isomorphism
\cite[Lemma~29.55.8, Tag~0AB1]{Stacks}.
Thus $\Gamma\simeq\PP^1_F$.

The restricted canonical bundle $\omega_R|_\Gamma$ has degree $-2$.
Indeed, differential transitivity gives
\[
 \Omega^1_{F/k}\otimes_F\OO_\Gamma\longrightarrow
 \Omega^1_{R/k}|_\Gamma\longrightarrow\Omega^1_{\Gamma/F}
 \longrightarrow0.
\]
The middle and final terms have ranks two and one and are locally
free. Their kernel is a line bundle, onto which the first, also a
line bundle, surjects. It is therefore isomorphic to that kernel.
Taking determinants proves the degree assertion.

For a point blowup, the local formula
$dx\wedge d(xv)=x\,dx\wedge dv$ gives an effective exceptional
canonical correction. Consequently
\[
 \rho^*\OO_X(A)\hookrightarrow
 \omega_R(-(m-2)\rho^*Y-\rho^*Z).
\]
A nonzero section in $W$ stays nonzero on the generic fibre, which
contains the generic point of $R$. The restriction of $\rho^*Z$ is
effective, including any horizontal components, so the target has
degree
\[
 -2-(m-2)-\deg(\rho^*Z|_\Gamma)<0.
\]
This contradicts the existence of that section. Hence the image of
$\psi$ has dimension two.

\paragraph{The square and the canonical system.}
Let $P=\psi^*\OO_B(1)$. Section ratios give an injection
$P\hookrightarrow\rho^*\OO_X(A)$ with effective exceptional zero
divisor $E_{\rm b}$; hence $\rho^*A=P+E_{\rm b}$.
Pullback is orthogonal to exceptional curves. Their intersection form
is negative definite, as the total transforms of successive exceptional
curves have diagonal matrix with entries $-1$.
The proper intersection-degree formula gives
\[
 P^2=[k(R):k(B)]\deg_{\OO_B(1)}B>0,
 \qquad A^2=P^2-E_{\rm b}^2>0
\]
\cite[Section~54.3, Tag~0AGP; Lemma~33.45.7, Tag~0BEL]{Stacks}.
No separability is needed. Products of the sections in $W$ have the
quadratic Hilbert growth of the homogeneous coordinate ring of $B$,
so $A$ is big.

For every $u\in H^0(X,\OO_X(Y))$, the one-form
$\beta(u)=(u/s)\omega$ is regular, since $m\geq2$.
The map $u\mapsto\beta(u)$ is injective. Closedness and $p\mid m$ give
\[
 d\beta(u)=s^{m-2}\gamma(u)=s^{m-2}s_Z\gamma'(u).
\]
Its kernel is $ks$, and multiplication by the fixed section does not
change the rational map. This proves the fixed-divisor, image and rank
assertions. Multiplication by $s^{m-p}s_Z$ places $W$ inside
$H^0(X,\omega_X(-(p-2)Y))$ without changing its rational map.
Adding effective divisors to $A$ proves the two bigness claims;
the canonical system contains $W_K$ and therefore has surface image
as well.
\end{proof}

The bigness consequence also has an argument that does not use
closedness or a distinguished divisor once a negatively twisted
cotangent section is known. It clarifies the role of the embedding.

\begin{lemma}\label{lem:negative-cotangent-adjoint}
Let $S$ be a smooth projective integral surface over an algebraically
closed field, let $H$ be very ample, let $a>0$ be an integer,
and let $Z$ be an effective Cartier divisor.
If $0\ne\alpha\in H^0(S,\Omega^1_S(-aH-Z))$, then
$|K_S+(2-a)H-Z|$ has a linear subsystem with two-dimensional image.
In particular $K_S+(2-a)H-Z$ is big. If $\alpha$ has no
divisorial zero, this divisor is also nef and has positive square.
\end{lemma}

\begin{proof}
Choose coordinate sections $s_0,\ldots,s_N$ for the embedding by $H$.
The expressions $\tau_{ij}=s_i\,ds_j-s_j\,ds_i$ glue to sections of
$\Omega^1_S(2H)$: the terms differentiating a change of frame cancel.
They generate this bundle, since on $s_i\ne0$ the $\tau_{ij}$ are,
in the induced trivialization, the differentials of the affine
coordinate functions $s_j/s_i$.
For the nonzero section $\alpha$, the wedges
$\alpha\wedge\tau_{ij}$ therefore give a rational map, regular where
$\alpha$ is nonzero, defined by a subsystem of $K_S+(2-a)H-Z$.

At such a point $x$, choose a nonzero tangent vector $v$ spanning the
kernel of $\alpha_x$. Up to one common scalar the coordinates of this
map are
\[
 s_i(x)\,ds_j(v)-s_j(x)\,ds_i(v).
\]
These are the Pl\"ucker coordinates of the projective line through
$x$ with tangent direction $v$. Thus each fibre of the rational map
lies in its represented line. The image cannot be a point, since
that would put a dense open of $S$ in one line.

If the image were a curve $B$ in the Grassmannian of lines, a general
fibre would be one-dimensional. Its closure would be the entire
represented line $C\subset S$. The universal family of lines is a
projective-line bundle over the Grassmannian; its restriction to a
dense smooth open of $B$ is flat and contained in $S\times B$.
The containment extends from the general fibres because the condition
that a line lie in $S$ is closed. These fibres are Cartier divisors on
the smooth surface, and the family is a relative Cartier divisor.
Indeed, flatness identifies the fibre of its ideal with the ideal of
the fibre; lifting a local generator and applying Nakayama's lemma
makes that ideal principal. Its generator is a nonzerodivisor in the
regular integral scheme $S\times B$ after shrinking to the smooth
base open. They form an algebraic family, so distinct members
$C,C'$ are numerically equivalent and
$C^2=C\cdot C'\geq0$.
For clarity, numerical equivalence can also be checked by restricting
the relative Cartier divisor to any fixed curve: the degree of its
line bundle is constant in this connected family.

Choose a general member $C$ that is not a component of $Z$ and meets
the open set where $\alpha\ne0$. There are only finitely many
components of $Z$, so this choice is possible.
Put $D=aH+Z$. Then $C\simeq\PP^1$, $H\cdot C=1$ and
$D\cdot C=a+Z\cdot C\geq a>0$. Its twisted conormal sequence is
\[
 0\longrightarrow\OO_C(-C-D)\longrightarrow
 \Omega^1_S(-D)|_C\longrightarrow\omega_C(-D)
 \longrightarrow0.
\]
The end terms have degrees $-C^2-D\cdot C<0$ and $-2-D\cdot C<0$.
Thus the middle term has no global section. This contradicts the
restriction of $\alpha$, which is nonzero at points of the original
fibre. The image therefore has dimension two. Its coordinate ring
embeds in the section ring of the adjoint and has quadratic Hilbert
growth, proving bigness. The argument uses no separability assumption.
If $\alpha$ has no divisorial zero, generation by the $\tau_{ij}$
shows that the wedge system has no fixed curve. The nefness and
exceptional-intersection argument in the preceding proof then gives
nefness and positive square. No closedness of $\alpha$ is required.
\end{proof}

\subsection{The structure of a conductor equality}
\label{subsec:conductor-equality-structure}

We now prove Theorem~\ref{thm:conductor-equality-structure}. The
positive square just obtained turns a determinant divisor into a
connected curve. Its constant functions split an Euler extension.
This supplies an ambient point through which all the tangent planes
pass, and makes the quotient and its possible singularity explicit.

The exclusion of $a=6$ will be completed in
Proposition~\ref{prop:conductor-six-excluded}; that argument uses
only the structural conclusions and $a\geq5$ proved here.

\begin{proof}[Proof of Theorem~\ref{thm:conductor-equality-structure}]
Put $H=\OO_X(Y)$, $m=pa$ and $d=Y^2$. Let $s_0$ define $Y$.
The equality conclusions of
Propositions~\ref{prop:normalized-mup-conductor}
and~\ref{prop:normalized-mup-singularity-defect} give
\[
 \beta=\omega/s_0^m\in H^0(X,\Omega^1_{X/k}\otimes H^{-m})
\]
whose tangential restriction to $Y$ is nowhere zero. Its zero scheme
$Z$ is finite and disjoint from $Y$: any divisorial zero would meet
the ample curve $Y$. Adjunction gives $K_X\cdot Y=(m-1)d$.
The annihilator of $\beta$ is the saturated invertible subsheaf
\[
 \mathcal F=H^m\otimes\omega_X^{-1}\longrightarrow T_X.
\]
Locally its generator is the Koszul relation between the two
coefficients of $\beta$, which have no common prime divisor.
It is a subbundle outside $Z$, and its normal projection on $Y$
is an isomorphism $\mathcal F|_Y\simeq H|_Y$.

Set
\[
 A=H\otimes\mathcal F^{-1}=\omega_X\otimes H^{-(m-1)},
 \qquad B=A\otimes H.
\]
Thus $A\cdot Y=0$ and $A|_Y\simeq\OO_Y$.
Proposition~\ref{prop:residual-adjoint}, with no additional
divisorial zero, shows that $B$ is nef and big.

\paragraph{Splitting the Euler extension.}
Let $V=H^0(X,H)$ and embed $X$ in $\PP(V^*)$. Pull back the
restricted Euler sequence along
$\mathcal F\to T_{\PP(V^*)}|_X$. This gives a rank-two bundle
$P$ and a commutative diagram
\[
\begin{array}{ccccccccc}
0&\longrightarrow&\OO_X&\longrightarrow&P&\longrightarrow&\mathcal F&\longrightarrow&0\\
&&\Vert&&\downarrow&&\downarrow\\
0&\longrightarrow&\OO_X&\longrightarrow&V^*\otimes H&\longrightarrow&T_{\PP(V^*)}|_X&\longrightarrow&0.
\end{array}
\]
The bundle $P$ remains locally free at $Z$, since it is an extension
of line bundles. Put $E=P^*\otimes H$. There are maps
\[
 0\longrightarrow A\longrightarrow E\longrightarrow H
 \longrightarrow0,\qquad V\otimes\OO_X\longrightarrow E,
\]
and their composite is evaluation. For $s\in V$, denote its
section of $E$ by $e_s$. Locally its components record the value
of $s$ and differentiation in the $\mathcal F$ direction. At a
zero of $s$, the derivative component is independent of the frame.
The section $e_0=e_{s_0}$ is nowhere zero: off $Y$ its value
component is nonzero, and on $Y$ its normal derivative is nonzero.
Consequently
\begin{equation}\label{eq:conductor-euler-extension}
 0\longrightarrow\OO_X\xrightarrow{e_0}E
 \xrightarrow{e_0\wedge-}B\longrightarrow0
\end{equation}
is exact. Its determinant sections $t_s=e_0\wedge e_s$ restrict
on $Y$ to the hyperplane sections $s|_Y$, under the identification
$B|_Y\simeq H|_Y$ given by $e_0$.

Choose $s$ so that $s|_Y$ has $d$ distinct simple zeros. Such
hyperplanes form a nonempty open set, since $Y$ is a smooth
embedded curve. The divisor $C=\operatorname{div}(t_s)$ is
nonempty and does not contain $Y$. Every component meets $Y$;
simple intersection with $Y$ therefore forces every component to
have multiplicity one. A Cartier divisor on a smooth surface has
no embedded components, so $C$ is reduced.

It is also connected. If $C=C_1+C_2$ were a disjoint decomposition
into nonzero effective divisors, then
$C_i^2=B\cdot C_i\geq0$. A zero square would make $C_i$
orthogonal to the positive-square class $B$, and Hodge index
would make $C_i$ numerically zero, contrary to $Y\cdot C_i>0$.
Two positive-square classes cannot be orthogonal either.
Thus $H^0(C,\OO_C)=k$.

On $C$, the image of $e_s$ in the quotient in
\eqref{eq:conductor-euler-extension} vanishes. Hence
$e_s|_C=\lambda e_0|_C$ for a scalar $\lambda\in k$.
Dividing $e_s-\lambda e_0$ by the equation $t_s$ gives a map
$B\to E$ whose composite with $E\to B$ is the identity.
Thus $E\simeq\OO_X\oplus B$.
The surjection $E\to H$ has a nonzero component $B\to H$;
otherwise its image would be generated by $s_0$ and would vanish
on $Y$. This gives a nonzero section of $A^{-1}$.
Its effective zero divisor has intersection zero with the ample
curve $Y$, so it is empty. Therefore
\begin{equation}\label{eq:conductor-linear-canonical}
 A\simeq\OO_X,\qquad B\simeq H,\qquad
 \mathcal F\simeq H,\qquad K_X\sim(m-1)Y.
\end{equation}
This proves linear equivalence, with no injectivity assumption on
restriction of line bundles to $Y$.

\paragraph{The constant tangent point.}
In $E\simeq\OO_X\oplus H$, the component $H\to H$ of the
surjection $E\to H$ is a nonzero scalar. It splits that surjection.
Dualizing and twisting therefore splits the original Euler pullback
compatibly. The inclusion $H\simeq\mathcal F\to T_{\PP(V^*)}|_X$
lifts to a map $H\to V^*\otimes H$, given by a constant nonzero
vector $v\in V^*$. Write $v$ also for its projective point.
The induced tangent direction at $x\ne v$ is that of the line
$xv$. In particular every embedded tangent plane of $X$ contains
$v$. This is the usual notion of a strange embedding
\cite[\S2.1 and Proposition~2.5, author pp.~2--5]{FurukawaStrange}.

The point $v$ lies outside the hyperplane defining $Y$, since the
direction on $Y$ is transverse to that hyperplane. Choose
coordinates with $v=[1:0:\cdots:0]$ and $Y=(x_0=0)\cap X$.
The corresponding section of $T_X\otimes H^{-1}$ has no zero
except possibly $v$. If $v\in X$, its zero ideal is exactly the
maximal ideal of $v$: the ambient coefficients $x_i/x_0$ generate
that ideal, and the tangent inclusion is locally a direct summand.
Thus $Z$ is empty or one reduced point. Let $\epsilon$ be its
length.

At this point choose two ambient affine coordinates $u,w$ forming
regular parameters of $X$. The radial field satisfies
$D(u)=-u$, $D(w)=-w$. The first jet of its annihilating form is
a nonzero scalar multiple of $u\,dw-w\,du$.
Since $v\notin Y$, the same is true of $\omega$.
The derivative of a local unit contributes zero at a zero of the
form. Thus $d\omega(v)=2c\,du\wedge dw$ with $c\ne0$.
Closedness proves
\begin{equation}\label{eq:conductor-epsilon-parity}
 \epsilon=1\quad\Longrightarrow\quad p=2.
\end{equation}
The zero-length formula gives, in both cases,
\begin{equation}\label{eq:conductor-c2X}
 \epsilon=c_2(\Omega_X^1\otimes H^{-m})=c_2(X)+md,
 \qquad c_2(X)=\epsilon-pa\,d.
\end{equation}

\paragraph{The quotient and the normalized torsor.}
Multiplying the twisted field by $x_0$ gives the restriction of
$\delta=x_0\partial/\partial x_0$, with $\delta^p=\delta$.
The homogeneous ideal of $X$ is stable under this derivation.
The eigenspace projections are polynomials in $\delta$, since
$T^p-T$ has distinct roots, so the ideal is graded by weights
modulo $p$. Hence the ambient $\mu_p$ action restricts to $X$.
Its quotient $q:X\to S$ is finite radicial of degree $p$ and
normal. On invariant affine rings, normality follows by
intersection with the invariant function field. The quotient is
projective; equivalently, relative Frobenius factors through it
and makes $S$ finite over $X^{(p)}$.

Let $K=k(X)$ and write the original torsor generically as
$t^p=u$, with $u\notin K^p$. Then $k(N)=L=K(u^{1/p})$.
The field $\delta$ annihilates $du/u$, and hence $u$.
Extending $u$ to a $p$-basis of $K$ shows that
\[
 K^\delta=K^p(u)=L^p.
\]
This identifies the schemes as well as their function fields.
For an affine open $\Spec A\subset S$, let its inverse image in
$N$ be $\Spec C$. Both rings are normal, $C$ is finite over $A$,
and $\operatorname{Frac}A=(\operatorname{Frac}C)^p$.
For $c\in C$, the element $c^p$ is integral over $A$ and lies
in its fraction field, so it belongs to $A$.
Conversely, every $a\in A$ has a $p$th root in
$\operatorname{Frac}C$, integral over $C$, and therefore in $C$.
Thus $A=C^p$. Since $k$ is perfect, these identifications give
\begin{equation}\label{eq:conductor-frobenius-quotient}
 S\simeq N^{(p)},\qquad q\pi=F_{N/k}.
\end{equation}

Away from $v$, the saturated foliation is regular, so its quotient
is smooth. The local result used here is Seshadri's regular
quotient theorem \cite[Proposition~6 and Theorem~2,
pp.~6-11--6-16]{SeshadriCartier}. In rank one, a normalized
generator has $\partial(t)=1$ and $\partial^p=0$; invariant
coordinates give a quotient with regular parameters $t^p,z$.
At $Y$ one can also use the multiplicative coordinates
$\delta(t)=t$, $\delta(z)=0$.

If $v\in X$, then $p=2$ and its completed invariant ring is
\[
 k[[u^2,uw,w^2]]\simeq k[[a,b,c]]/(ac-b^2).
\]
Indeed, invariant monomials have even total degree, and are
products of $u^2,uw,w^2$, with the displayed sole relation.
Invariants commute with completion through the weight-zero
direct summand. This is an $A_1$ singularity. Blowing up $v$
in $X$, with exceptional curve $R$, resolves the quotient:
on the chart $w=uz$ the lifted field satisfies
$\delta(u)=-u$, $\delta(z)=0$, and the quotient chart has
coordinates $u^2,z$. The other chart is analogous.
Let $q':X'\to T'$ be this smooth quotient and $\rho:T'\to S$
the induced resolution. Its exceptional curve $E$ satisfies
$q'^*E=2R$ and $E^2=-2$.
It is the minimal resolution. It is crepant as well: since $S$
is Gorenstein, write $K_{T'}=\rho^*K_S+cE$; adjunction gives
$K_{T'}\cdot E=0$, whence $c=0$.
For $\epsilon=0$ set $X'=X$ and $T'=S$.
Equation~\eqref{eq:conductor-frobenius-quotient} transfers these
smoothness and singularity statements to $N$; its minimal
resolution $T$ has Frobenius twist $T'$.

\paragraph{Canonical class and Chern numbers.}
The normalized conductor formula gives
\[
 K_N=\pi^*(K_X-a(p-1)Y)\sim(a-1)\pi^*Y.
\]
For $a\geq2$ it is ample. The crepant resolution has nef and
big canonical divisor, so $T$ and $T'$ are minimal surfaces of
general type. Their canonical squares are
\[
 K_T^2=K_{T'}^2=K_N^2=p(a-1)^2d.
\]

We compute their second Chern numbers directly on smooth surfaces.
For a regular rank-one foliation $F\subset T_{X'}$, let $Q$
be its annihilator in $\Omega^1_{X'}$.
The invariant coordinates just used give exact sequences
\[
 0\to Q\to\Omega^1_{X'}\to F^{-1}\to0,
 \qquad
 0\to F^{-p}\to q'^*\Omega^1_{T'}\to Q\to0.
\]
Here $Q=\omega_{X'}\otimes F$; the $p$th power in the second
kernel follows from the transition of the local coordinate $t^p$.
Writing $f=c_1(F)$, determinants and second Chern classes give
\[
 q'^*K_{T'}=K_{X'}-(p-1)f,\qquad
 c_2(X')=-f\cdot(K_{X'}+f),\qquad c_2(T')=c_2(X').
\]
For $\epsilon=0$, $f=Y$. For $\epsilon=1$, the lifted
field has fixed divisor $b^*Y+R$, so $f=b^*Y+R$ and
$K_{X'}=b^*K_X+R$.
Using $R^2=-1$ and $b^*Y\cdot R=0$, one obtains
\begin{equation}\label{eq:conductor-structure-invariants}
\begin{gathered}
 K_T^2=p(a-1)^2d,\qquad c_2(T)=-pa\,d+2\epsilon,\\
 12\chi(\OO_T)=p(a^2-3a+1)d+2\epsilon.
\end{gathered}
\end{equation}
The last equality is Noether's formula.
As a separate check, Noether's formula on $X$ and
Corollary~\ref{cor:normalized-equality-euler} give
$12\chi(\OO_N)=p(a^2-3a+1)d+p\epsilon$;
equation~\eqref{eq:conductor-epsilon-parity} makes the two
expressions equal.

\paragraph{The first four conductors are impossible.}
If $a=1$, $K_T$ is numerically trivial.
The classification of smooth projective surfaces with numerically
trivial canonical divisor gives $c_2(T)\geq0$
\cite[\S2.1, Table~1, author p.~3]{MatsumotoMuN2023}.
The table has Betti pairs $(b_1,b_2)=(4,6),(0,22),(0,10),(2,2)$,
giving Euler numbers $0,24,12,0$.
For $\epsilon=0$, this contradicts $c_2(T)=-pd<0$.
For $\epsilon=1$, $p=2$ and $2-2d\geq0$, hence $d=1$.
But $H|_Y$ is very ample of degree one, so $Y$ is a line.
This contradicts $2g(Y)-2=pa\,d>0$.

For $a\geq2$, Gu--Sun--Zhou's inequality applies to the smooth
minimal surface $T$ of general type:
\[
 K_T^2\leq32\chi(\OO_T)
 \qquad\text{\cite[Theorem~1.5, p.~614;
 Corollary~3.4, pp.~623--624]{GuSunZhou}.}
\]
In particular $\chi(\OO_T)>0$.
If $a=2$, equation~\eqref{eq:conductor-structure-invariants}
gives $12\chi(\OO_T)=-pd+2\epsilon\leq0$, a contradiction.
If $a=3$ and $\epsilon=0$, it gives
$K_T^2=4pd>32pd/12=32\chi(\OO_T)$.
If $a=3$ and $\epsilon=1$, it gives
\[
 K_T^2=8d,\qquad \chi(\OO_T)=(d+1)/6.
\]
The inequality forces $d\leq2$, but neither positive integer
$d=1,2$ makes this Euler characteristic integral.
Thus $a\geq4$, proving ampleness and minimality.

We have established $a\geq4$. Suppose $a=4$, and let
$\rho:T\to N$ be the minimal resolution, with $\rho$ the identity
when $N$ is smooth. Its canonical identity and crepancy give
\[
 K_T\sim3A,\qquad A=\rho^*\pi^*Y,
\]
where $A$ is an integral Cartier divisor class. The invariant
formulas give
\[
 K_T^2=9pd,\qquad 12\chi(\OO_T)=5pd+2\epsilon.
\]
Since $\epsilon=1$ implies $p=2$, it follows that
\[
 12\bigl(K_T^2-18\chi(\OO_T)\bigr)
 =18pd-36\epsilon=18p(d-\epsilon)>0.
\]
For $\epsilon=0$, positivity follows from $d>0$.
For $\epsilon=1$, the value $d=1$ would make the very ample
curve $Y$ a line, contradicting $2g(Y)-2=pa\,d>0$.

The surface $T$ is smooth minimal of general type.
Gu--Sun--Zhou's theorem therefore gives a genus-two Albanese
fibration
\cite[Theorem~1.5, p.~614; Corollary~3.4 and its proof,
pp.~623--624]{GuSunZhou}.
For a fibre divisor $F$, one has $F^2=0$ and
$K_T\cdot F=2$ by adjunction.
The genus is the arithmetic genus in the flat-fibration
convention of that theorem, so this calculation does not require
a smooth fibre. It contradicts $K_T\sim3A$, since
$A\cdot F$ is an integer.
Therefore $a\geq5$.

\paragraph{The divisibility of the degree.}
If $v\notin X$, projection from $v$ is a morphism defined by
$x_1,\ldots,x_r$, with pullback of $\OO(1)$ equal to $H$.
Its coordinate ratios are invariant under $\mu_p$, so it factors
through $q:X\to S$. For the descended line bundle $L$ on $S$,
$q^*L=H$. Hence $d=pL^2$.
If $v\in X$, the base ideal of the projection is the maximal
ideal of that reduced point. On $X'$ the resolved projection
has line bundle $b^*H-R$ and factors through $q':X'\to T'$.
Its descended line bundle satisfies $q'^*L=b^*H-R$; thus
$d-1=pL^2$. This proves $p\mid(d-\epsilon)$.
No very ampleness of the descended line bundle is required.
\end{proof}

\begin{corollary}\label{cor:conductor-not-complete-intersection}
In Theorem~\ref{thm:conductor-equality-structure}, the embedding
of $X$ by $\OO_X(Y)$ is not a complete intersection.
\end{corollary}

\begin{proof}
The embedding is strange, as the proof shows, and is not linear
since $K_X\cdot Y=(pa-1)d>0$.
Furukawa's theorem states that a nonlinear smooth complete
intersection is strange exactly when it is an odd-dimensional
quadric in characteristic two
\cite[Theorem~1.1, author p.~2, proved in \S3]{FurukawaStrange}.
Our surface has even dimension, so it cannot be such a complete
intersection.
\end{proof}


\begin{corollary}\label{prop:higher-conductor-positive-genus-exclusion}
Let $X$ be a smooth projective integral surface over an algebraically
closed field of characteristic $p>0$, let $Y\subset X$ be a smooth
integral very ample divisor, and let $M\ne\OO_X$ have order $p$ with
$M|_Y\simeq\OO_Y$. Let $pa$ be the exact vanishing order of its
logarithmic form along $Y$, and let $N$ be the normalization of its
$\mu_p$-torsor. If $2g(Y)-2=paY^2$, then $N$ is not birationally ruled.
\end{corollary}

\begin{proof}
Theorem~\ref{thm:conductor-equality-structure} makes the minimal
resolution of $N$ a surface of general type. A birationally ruled
surface has Kodaira dimension $-\infty$, so it cannot occur.
\end{proof}

\begin{corollary}\label{lem:divided-differential-subsystem}
In Proposition~\ref{prop:residual-adjoint}, if $m=p$ and the tangential
restriction of $\omega/s^p$ to $\omega_Y(-pY)$ is nowhere vanishing,
then $Z=0$. The residual system is a subsystem of
$\mathcal L=\omega_X(-(p-2)Y)$, has no base point on $Y$, and restricts
to its given embedding. In particular $\mathcal L$ is nef and big and
\begin{equation}\label{eq:divided-subsystem-positive-square}
 \mathcal L^2>0.
\end{equation}
\end{corollary}

\begin{proof}
A nowhere-vanishing tangential section makes
$2g(Y)-2-pY^2=0$. Proposition~\ref{prop:residual-adjoint} gives
$Z\cdot Y=0$. Since $Y$ is ample and $Z$ effective, $Z=0$.
The restriction formula in its proof has a nowhere-vanishing common
factor, so there is no base point on $Y$. All the other conclusions
follow from that proposition.
\end{proof}

\begin{proof}[Proof of Theorem~\ref{thm:strict-picard-all}]
The nontrivial order-$p$ class gives the nonzero closed logarithmic
form of Proposition~\ref{prop:normalized-mup-conductor}. Its full
restriction to $Y$ vanishes, by the conormal sequence and negative
conormal degree. Proposition~\ref{prop:residual-adjoint} proves the
adjoint and canonical-image assertions.
The conductor inequality is $2g(Y)-2\geq paY^2\geq pY^2$.
If the last bound were an equality, then $a=1$ and equality would
hold with the exact conductor. This contradicts
Theorem~\ref{thm:conductor-equality-structure}, which requires
$a\geq5$. Hence $2g(Y)-2>pY^2$.
\end{proof}
